% Caso de Uso: Agregar Desparasitación
% Sistema: Hotel Perros

%-------------------------------------- COMIENZA descripción del caso de uso.

\begin{UseCase}{CU10}{Agregar Desparasitación}{
    Permitir al Propietario registrar un tratamiento de desparasitación realizado a su mascota para mantener un historial médico completo y actualizado.
}
    \UCitem{Versión}{\color{Gray}1.0}
    \UCitem{Autor}{\color{Gray}Vazquez Villagran Jorge}
    \UCitem{Supervisa}{\color{Gray}Castillo Aldaba Jared}
    \UCitem{Actor}{\hyperlink{Propietario}{Propietario}}
    \UCitem{Propósito}{Registrar información sobre tratamientos de desparasitación aplicados a las mascotas para llevar un control médico adecuado.}
    \UCitem{Entradas}{Tipo de desparasitación, Producto utilizado, Fecha de aplicación, Próxima fecha programada.}
    \UCitem{Origen}{Teclado}
    \UCitem{Salidas}{Confirmación de registro exitoso, comprobante de desparasitación registrada, historial actualizado de la mascota.}
    \UCitem{Destino}{Pantalla e impresora para comprobante}
    \UCitem{Precondiciones}{
        \begin{itemize}
            \item El propietario debe estar autenticado en el sistema.
            \item El propietario debe tener al menos una mascota registrada.
            \item La mascota debe pertenecer al propietario autenticado.
        \end{itemize}
    }
    \UCitem{Postcondiciones}{
        \begin{itemize}
            \item La desparasitación queda registrada en el historial médico de la mascota.
            \item Se actualiza la fecha de próxima desparasitación recomendada.
            \item El sistema puede enviar recordatorios basándose en la próxima fecha programada.
        \end{itemize}
    }
    \UCitem{Errores}{
        \begin{itemize}
            \item MSG1: La mascota no existe o no pertenece al propietario.
            \item MSG2: Formato de fecha inválido.
            \item MSG3: La próxima fecha debe ser posterior a la fecha de aplicación.
            \item MSG4: Campos requeridos faltantes.
        \end{itemize}
    }
    \UCitem{Tipo}{Caso de uso primario}
    \UCitem{Observaciones}{
        \begin{itemize}
            \item El tipo de desparasitación puede ser: interna, externa o ambas.
            \item Se recomienda registrar el producto comercial utilizado para futuras referencias.
            \item El sistema puede sugerir la próxima fecha basándose en recomendaciones veterinarias estándar.
        \end{itemize}
    }
\end{UseCase}

%--------------------------------------
\begin{UCtrayectoria}
    \UCpaso[\UCactor] Accede al sistema con su correo electrónico y contraseña mediante la \IUref{IU01}{Pantalla de Inicio de Sesión}\label{CU08Login}.
    \UCpaso[\UCactor] Confirma el acceso presionando el botón \IUbutton{Iniciar Sesión}.
    \UCpaso Verifica las credenciales del propietario con base en la regla \BRref{BR01}{Autenticar Usuario}.
    \UCpaso Despliega la \IUref{IU10}{Pantalla Principal del Propietario} con el listado de sus mascotas registradas.
    \UCpaso[\UCactor] Selecciona la mascota a la que desea agregar el registro de desparasitación \Trayref{A}\label{CU08SeleccionarMascota}.
    \UCpaso Verifica que la mascota pertenezca al propietario autenticado con base en la regla \BRref{BR10}{Verificar Propiedad de Mascota} \Trayref{B}.
    \UCpaso Despliega la \IUref{IU25}{Pantalla de Detalle de Mascota} con las opciones del historial médico.
    \UCpaso[\UCactor] Selecciona la opción ``Agregar Desparasitación'' en el menú de historial médico.
    \UCpaso Despliega el \IUref{IU26}{Formulario de Registro de Desparasitación}.
    \UCpaso[\UCactor] Ingresa el tipo de desparasitación (interna, externa o ambas) \Trayref{C}.
    \UCpaso[\UCactor] Ingresa el nombre del producto utilizado (opcional).
    \UCpaso[\UCactor] Selecciona la fecha de aplicación del tratamiento \Trayref{D}.
    \UCpaso[\UCactor] Selecciona la próxima fecha programada para desparasitación (opcional) \Trayref{E}.
    \UCpaso[\UCactor] Confirma el registro presionando el botón \IUbutton{Guardar Desparasitación}.
    \UCpaso Valida que todos los campos requeridos estén completos con base en la regla \BRref{BR30}{Validar Datos de Desparasitación} \Trayref{F}.
    \UCpaso Valida que las fechas ingresadas sean coherentes con base en la regla \BRref{BR31}{Validar Fechas de Desparasitación} \Trayref{G}.
    \UCpaso Registra la desparasitación en la base de datos asociada a la mascota.
    \UCpaso Muestra el mensaje ``Desparasitación registrada exitosamente'' mediante la \IUref{IU27}{Pantalla de Confirmación de Registro}.
    \UCpaso Actualiza el historial médico de la mascota en la \IUref{IU25}{Pantalla de Detalle de Mascota}.
    \UCpaso Pregunta al propietario si desea imprimir un comprobante del registro.
    \UCpaso[\UCactor] Indica que desea imprimir el comprobante \Trayref{H}.
    \UCpaso Genera e imprime el \IUref{IU28}{Comprobante de Desparasitación Registrada} con base en la regla \BRref{BR32}{Formato de Comprobante de Desparasitación}.
\end{UCtrayectoria}

%--------------------------------------
\begin{UCtrayectoriaA}{A}{El propietario no tiene mascotas registradas}
    \UCpaso Muestra el Mensaje {\bf MSG5-}``No tiene mascotas registradas en el sistema. Debe registrar al menos una mascota para continuar.''.
    \UCpaso Despliega la opción de \IUbutton{Registrar Nueva Mascota}.
    \UCpaso[\UCactor] Selecciona registrar una nueva mascota o cancelar la operación.
    \UCpaso Si selecciona registrar, redirige al \UCref{CU02}{Registrar Mascota}.
    \UCpaso Si cancela, continúa en el paso \ref{CU08Login} del \UCref{CU08}.
\end{UCtrayectoriaA}

%--------------------------------------
\begin{UCtrayectoriaA}{B}{La mascota no pertenece al propietario}
    \UCpaso Muestra el Mensaje {\bf MSG1-}``La mascota seleccionada no existe o no pertenece a su cuenta.''.
    \UCpaso[\UCactor] Oprime el botón \IUbutton{Aceptar}.
    \UCpaso Continúa en el paso \ref{CU08SeleccionarMascota} del \UCref{CU08}.
\end{UCtrayectoriaA}

%--------------------------------------
\begin{UCtrayectoriaA}{C}{El propietario decide no especificar el tipo}
    \UCpaso[\UCactor] Deja el campo de tipo vacío o selecciona ``No especificado''.
    \UCpaso El sistema asigna el valor por defecto ``General''.
    \UCpaso Continúa con el paso siguiente del \UCref{CU08}.
\end{UCtrayectoriaA}

%--------------------------------------
\begin{UCtrayectoriaA}{D}{Fecha de aplicación inválida}
    \UCpaso El sistema detecta que la fecha ingresada no tiene un formato válido.
    \UCpaso Muestra el Mensaje {\bf MSG2-}``La fecha de aplicación ingresada no es válida. Por favor use el formato DD/MM/AAAA.''.
    \UCpaso Resalta el campo de fecha en color rojo.
    \UCpaso[\UCactor] Corrige la fecha de aplicación.
    \UCpaso Continúa con el paso siguiente del \UCref{CU08}.
\end{UCtrayectoriaA}

%--------------------------------------
\begin{UCtrayectoriaA}{E}{Próxima fecha anterior a la fecha de aplicación}
    \UCpaso El sistema detecta que la próxima fecha programada es anterior o igual a la fecha de aplicación actual.
    \UCpaso Muestra el Mensaje {\bf MSG3-}``La próxima fecha de desparasitación debe ser posterior a la fecha de aplicación.''.
    \UCpaso Resalta los campos de fecha en color rojo.
    \UCpaso[\UCactor] Corrige la próxima fecha programada.
    \UCpaso Continúa con el paso siguiente del \UCref{CU08}.
\end{UCtrayectoriaA}

%--------------------------------------
\begin{UCtrayectoriaA}{F}{Campos requeridos faltantes}
    \UCpaso El sistema detecta que faltan campos obligatorios por completar.
    \UCpaso Muestra el Mensaje {\bf MSG4-}``Debe completar todos los campos requeridos: [{\em lista de campos faltantes}].''.
    \UCpaso Resalta los campos faltantes en color rojo.
    \UCpaso[\UCactor] Completa los campos requeridos.
    \UCpaso Continúa con el paso siguiente del \UCref{CU08}.
\end{UCtrayectoriaA}

%--------------------------------------
\begin{UCtrayectoriaA}{G}{Error en validación de fechas}
    \UCpaso El sistema detecta inconsistencias en las fechas (ejemplo: fecha futura para aplicación).
    \UCpaso Muestra el Mensaje {\bf MSG6-}``La fecha de aplicación no puede ser una fecha futura.''.
    \UCpaso[\UCactor] Corrige la fecha de aplicación.
    \UCpaso Continúa con el paso siguiente del \UCref{CU08}.
\end{UCtrayectoriaA}

%--------------------------------------
\begin{UCtrayectoriaA}{H}{El propietario no desea imprimir comprobante}
    \UCpaso[\UCactor] Indica que no desea imprimir el comprobante.
    \UCpaso Muestra mensaje ``El comprobante está disponible en el historial de la mascota para consulta posterior.''.
    \UCpaso[\UCactor] Oprime el botón \IUbutton{Aceptar}.
    \UCpaso Termina el caso de uso.
\end{UCtrayectoriaA}

%--------------------------------------
% Puntos de extensión
\subsection{Puntos de extensión}

\UCExtenssionPoint{
    % Cuando:
    Desea consultar el historial completo de desparasitaciones de la mascota.
}{
    % Durante la región:
    Del paso 7 al paso 19.
}{
    % Casos de uso a los que extiende:
    \hyperlink{CU09}{CU09 Consultar Historial de Desparasitaciones}.
}

\UCExtenssionPoint{
    % Cuando:
    Desea editar o eliminar una desparasitación registrada previamente.
}{
    % Durante la región:
    Después del paso 19.
}{
    % Casos de uso a los que extiende:
    \hyperlink{CU10}{CU10 Modificar Desparasitación}, \\
    \hyperlink{CU11}{CU11 Eliminar Desparasitación}.
}

\UCExtenssionPoint{
    % Cuando:
    Desea configurar recordatorios automáticos para la próxima desparasitación.
}{
    % Durante la región:
    Después del paso 22.
}{
    % Casos de uso a los que extiende:
    \hyperlink{CU45}{CU45 Configurar Recordatorios Médicos}.
}

%-------------------------------------- TERMINA descripción del caso de uso.